% Title: Local Grocery Store Curry
% Description: This is defined from my head from 2020-2025. Vegetarian friendly, can be gluten free with commonplace modifications to starch binder for roux.
% Author: Wesley Jones
% Date: Oct. 2025
\documentclass{article}


\usepackage[nonumber,noindex,nocontents]{cuisine}
\usepackage[margin=1in]{geometry}
\usepackage{endnotes}

% Set widths
\RecipeWidths  
  {1\textwidth}    %{Total recipe width}             
  {0.1\textwidth}  %{Number of servings width}                   
  {0.05\textwidth} %{Step number width}             
  {0.15\textwidth} %{Ingredient width} 
  {0.1\textwidth}  %{Quantity width} 
  {0.1\textwidth}  %{Units width}


% End notes instead of footnotes
\let\footnote=\endnote


\begin{document}
\begin{recipe}{Local Grocery Store Curry}{\textit{Serves:  4}}{\textit{Preparation Time: 40 minutes}}
  
  \freeform
    Many curries can be made with a large variety of ingredients.
    This is a simple base to get started with --- all ingredients are intended to be readily available at your local grocery store.

  \Ingredient{2 cups chickpeas, drained, peeled, \& rinsed}
    
    Apply cooking oil to sheet pan.
    Place chickpeas on pan and toss them around to coat in oil.
    Roast chickpeans in 375\0 oven for 12 minutes.
    At minute 6, remove and apply salt and pepper, rotating peas to coat again in oil, salt, and pepper.
    Remove and set aside when done.

  \Ingredient{1 block firm tofu, pressed
    \endnote{\textbf{Substitute:} one half to 1 lbs. your choice of meat.}
  }
  \Ingredient{1 dash chili powder, or other warm spice}
  
    Cube into \fr{1}{2} inch or smaller peices and cook the tofu in oil on a medium skillet.
    Season with salt, pepper, and a dash of chili powder (or alternative warm spice).
    It will cook best when allowed to sit untended in the skillet for 5 minute increments before flipping peices.
    Tofu is done when it has a golden brown crust on each side.
    Remove from skillet and set aside when done.
  
  \Ingredient{1-2 cups chopped vegetables 
    \endnote{
      \textbf{Suggestions for  vegetables:} onions, green bell peppers, red bell peppers, and carrots. 
        Additional options can include nearly any vegetables that need to be eaten.
    }
  }  
  \Ingredient{2 cloves garlic, chopped or minced 
    \endnote{
      \textbf{Substitute:} clove garlic with garlic powder.
      Use two to three tsp., and add when returning tofu and chickpeas to curry.
    }
  }
  \Ingredient{1 Tbsp curry powder}
  \Ingredient{\fr{1}{2} tsp chili powder}
  \Ingredient{\fr{1}{4} tsp cumin}
    
    Sautee vegetables in skillet when tofu is done cooking at medium heat, with some additional oil or reserved oil from previous cook.
    Once vegetables begin to soften you can add salt and pepper.
    Additionally, apply more heat and leave vegetables undisturbed to add a slight char (this improves their flavor).
    Add garlic, stir.
    Create a gap in the vegetables and add the curry powder, chili powder, and cumin.
    Allow these to ``bloom'' for a minute before stirring into the vegetables.

  \Ingredient{1 cup milk}
  \Ingredient{2 Tbsp butter}
  \Ingredient{2 Tbsp flour}
    Add butter to skillet with vegetables. Allow butter to melt and foam. Reduce
    heat to medium low. Once butter melts and foaming resides, add flour and
    stir. The vegetables and flour will become clumpy, this is fine. Once flour
    is coat and combined with melted butter, begin to add milk in, slowly. Stir
    a small amount of milk in and stir until it is gone. Repeat this, adding
    a small amount and stiring until you have the start of a base (a roux). You
    will begin to see a smooth white sauce develop. Once you've added half the
    milk and have a sauce started you can dump the rest. Allow the sauce to
    thicken slightly for a few minutes.
  \Ingredient{\fr{1}{2} cup tomato sauce}
  \Ingredient{\fr{1}{4} cup sour cream}
  \Ingredient{1 cup broth or stock
    \endnote{chicken, beef, or vegetable are  all fine}
  }
  \Ingredient{1 Tbsp parsley}
    Add tomato sauce, the amount depends on your preferences. Sour cream should
    be added immediately with cold or room temperature tomato sauce to prevent
    curdling. Sour cream can work to adjust flavor preferences, add in
    ``dollops'', with two big ones being a \fr{1}{4} cup. Add a small amount of
    broth or stock and stir all ingredients into sauce. Return tofu (or prepared
    meat) and chickpeas to skillet. Stir. As ingredients warm and thicken you can
    add stock or both to attain a desired consistency. Add parsley.
    \begin{itemize}
      \item{\textit{For serving atop rice:} looser is nice to create stew-like dish.
          Additionally, add extra salt and pepper to account for unsalted
        rice.}
      \item{\textit{For serving with naan:} a thicker consistency might be beneficial.
        Again, attend to seasoning finishes to match provided medium.}
    \end{itemize}
\end{recipe}

\theendnotes

\end{document}
